\begin{kbvrule}
\noindent\kbvbullet\hspace{0.5em}ವಾಕ್ಯವು ಕರ್ತೃ, ಕರ್ಮ, ಕ್ರಿಯಾಪದಗಳನ್ನೊಳಗೊಂಡ ಪದ ಸಮುದಾಯವೆಂಬುದನ್ನು ತಿಳಿಮುದಾಯಿತು. ಆದರೆ ಕೆಲವು ವಾಕ್ಯಗಳಲ್ಲಿ ಕರ್ಮಪದದ ಅಗತ್ಯವೇ ಇರುವುದಿಲ್ಲ. ಉದಾಹರಣೆಗೆ, ಗಿಣಿಯು ಹಾರುತ್ತಿದೆ. ಕತ್ತೆ ಒದೆಯಿತು. ರಾಮನು ಬರುತ್ತಾನೆ. ವಿಮಾನ ಇಳಿಯಿತು. ನಾಯಿಯು ಓಡಿತು. ಹೂವು ಅರಳಿತು.
\end{kbvrule}

\indent ಈ ವಾಕ್ಯಗಳಲ್ಲಿ ಕರ್ಮಪದಗಳಿಲ್ಲ ; ಇಲ್ಲಿರುವುದು ಕರ್ತೃ ಮತ್ತು ಕ್ರಿಯಾಪದಗಳು ಮಾತ್ರ... ಆದರೂ ಈ ವಾಕ್ಯಗಳು ಅರ್ಥಪೂರ್ಣವಾಗಿವೆ. {[}ಈ ಬಗ್ಗೆ ಹೆಚ್ಚಿನ ವಿವರಗಳನ್ನು ಮುಂದಿನ ಪಾಠದಲ್ಲಿ ತಿಳಿಯಬಹುದು.{]}\par

\kbvexercise{ಅಭ್ಯಾಸ ಪತ್ರಿಕೆ ೪}

\vspace{2.5pt plus 3pt minus 0.5pt}\noindent\hspace*{1.2em}\kbvnum{(1)}\hspace{0.55em}ಈ ವಾಕ್ಯಗಳಲ್ಲಿರುವ ಕರ್ತೃಪದ, ಕ್ರಿಯಾಪದ ಮತ್ತು ಕರ್ಮ\par

\kbvsechead{ಪದಗಳನ್ನು ಗುರ್ತಿಸಿ :}

\vspace{2.5pt plus 3pt minus 0.5pt}\noindent\hspace*{1.2em}\kbvnum{1.}\hspace{0.55em}ಚಿರತೆಯು ಹಸುವನ್ನು ಕೊಂದಿತು.\par

\vspace{2.5pt plus 3pt minus 0.5pt}\noindent\hspace*{1.2em}\kbvnum{2.}\hspace{0.55em}ಭೀಮಸೇನನು ಹಿಡಂಬಿಯನ್ನು ಮದುವೆಯಾದನು.\par

\vspace{2.5pt plus 3pt minus 0.5pt}\noindent\hspace*{1.2em}\kbvnum{3.}\hspace{0.55em}ಹೂವು ಅರಳಿತು ; ಸುಗಂಧ ಹರಡಿತು.\par

\vspace{2.5pt plus 3pt minus 0.5pt}\noindent\hspace*{1.2em}\kbvnum{4.}\hspace{0.55em}ರೈಲುಗಾಡಿ ಚಲಿಸಿತು.\par

\vspace{2.5pt plus 3pt minus 0.5pt}\noindent\hspace*{1.2em}\kbvnum{5.}\hspace{0.55em}ಕಳ್ಳನು ಸರವನ್ನು ಕದ್ದನು.\par

\vspace{2.5pt plus 3pt minus 0.5pt}\noindent\hspace*{1.2em}\kbvnum{6.}\hspace{0.55em}ಬಟ್ಟೆಯನ್ನು ಕತ್ತರಿಯಿಂದ ಕತ್ತರಿಸುತ್ತೇವೆ.\par

\vspace{2.5pt plus 3pt minus 0.5pt}\noindent\hspace*{1.2em}\kbvnum{7.}\hspace{0.55em}ಸಕ್ಕರೆಯಿಂದ ಮಿಠಾಯಿ ತಯಾರಿಸುತ್ತಾರೆ.\par

\vspace{2.5pt plus 3pt minus 0.5pt}\noindent\hspace*{1.2em}\kbvnum{8.}\hspace{0.55em}ಅರ್ಚಕನು ಹೂವುಗಳಿಂದ ದೇವರನ್ನು ಪೂಜಿಸಿದನು.\par

\vspace{2.5pt plus 3pt minus 0.5pt}\noindent\hspace*{1.2em}\kbvnum{9.}\hspace{0.55em}ಹಡಗು ನೀರಿನಲ್ಲಿ ಮುಳುಗಿತು.\par

\vspace{2.5pt plus 3pt minus 0.5pt}\noindent\hspace*{1.2em}\kbvnum{10.}\hspace{0.55em}ಗಾಳಿಯು ವೇಗವಾಗಿ ಬೀಸುತ್ತಿದೆ.\par

\vspace{2.5pt plus 3pt minus 0.5pt}\noindent\hspace*{1.2em}\kbvnum{(2)}\hspace{0.55em}ಕೆಳಗೆ ಕೊಟ್ಟಿರುವ ಕ್ರಿಯಾಪದಗಳಿಗೆ ಸರಿಹೊಂದುವ ಕರ್ತೃ,\par

\kbvsechead{ಕರ್ಮ ಪದಗಳನ್ನು ಆಯ್ಕೆಮಾಡಿ ವಾಕ್ಯಗಳನ್ನು ಬರೆಯಿರಿ :}

\vspace{4pt plus 4pt minus 1pt}\noindent{\kbvsecfont\fontsize{11pt}{14.5pt}\selectfont\color{kbvaccentdark} 1.\hspace{0.45em}ಕಚ್ಚಿತು (ಮಾದರಿ :}\hspace{0.5em}ಮುಂಗುಸಿಯು ಹಾವನ್ನು ಕಚ್ಚಿತು.)\par
